\section{偏差---方差、正则化与模型选择}
\label{sec:11-Mathematical-Statistics/06-Resampling-and-Model-Selection/04}

上一节那条流水线已经关好了每一道门，只剩最后一个旋钮：ridge 的正则强度 \(\lambda\)。你铺了 \(20\) 个候选值，每个跑一遍 \(5\) 折交叉验证，得到一条曲线。最低点落在 \(\lambda=0.03\)，AUC \(0.802\)；隔壁 \(\lambda=0.1\) 是 \(0.799\)。差 \(0.003\)，你选了前者。

一个季度后回看，\(\lambda=0.1\) 那一版在线上稳定高出一小截。这不奇怪：\(0.003\) 的差距远小于交叉验证估计自身的波动，你挑的与其说是更好的 \(\lambda\)，不如说是在这组折上运气更好的 \(\lambda\)。

这一节把前面几节的工具收在一处。旋钮为什么会有一个中间的最优位置，交叉验证选中的那个点在什么意义上是最优的，以及在同一份验证数据上反复挑选时，这个``最优''会怎样自己往下沉。

\subsection{误差被拆成三块，其中两块朝相反方向动}

设数据由 \(Y=f(X)+\varepsilon\) 生成，\(\E[\varepsilon\given X]=0\)，\(\Var(\varepsilon\given X=x_0)=\sigma^2\)。训练集是随机的，因而学出来的预测器 \(\hat f\) 也是随机的。在一个固定的新点 \(x_0\) 上，取新响应 \(Y_0=f(x_0)+\varepsilon_0\)，把误差拆成三段再展开平方，

\[
\begin{aligned}
Y_0-\hat f(x_0)
&=\underbrace{\bigl(Y_0-f(x_0)\bigr)}_{\varepsilon_0}
+\bigl(f(x_0)-\E[\hat f(x_0)]\bigr)
+\bigl(\E[\hat f(x_0)]-\hat f(x_0)\bigr),\\
\E\bigl[(Y_0-\hat f(x_0))^2\bigr]
&=\sigma^2+\bigl(\E[\hat f(x_0)]-f(x_0)\bigr)^2+\Var\bigl(\hat f(x_0)\bigr).
\end{aligned}
\]

对一遍账。含 \(\varepsilon_0\) 的两个交叉项为零，用的是新响应的噪声与训练集独立且条件均值为零——如果验证数据与训练数据之间有泄漏，这一条就不成立，上一节整节都在保护它。第三段与其均值的交叉项为零，用的只是方差的定义。期望是对训练集的随机性和新响应的噪声同时取的，所以``偏差''和``方差''说的是重复采集训练数据时预测器的平均位置与摆动幅度，而不是这一份数据上的任何可观测量。

三项里 \(\sigma^2\) 与你选什么模型无关，是数据本身的不可约噪声；把模型调复杂，平方偏差下降而方差上升，两者反向。总误差因此常呈现一个中间的最低点。

\subsection{那条 U 形曲线上有两个位置容易被混淆}

\begin{center}
\begin{tikzpicture}[x=1.15cm,y=1.05cm,font=\small,>=stealth]
  \draw[->,black!55] (0,0) -- (7.2,0);
  \draw[->,black!55] (0,0) -- (0,3.35) node[above,font=\footnotesize] {误差};
  \node[font=\footnotesize,black!70] at (3.5,-0.4) {模型复杂度（正则强度减弱 \(\rightarrow\)）};
  % 真实风险
  \draw[dashed,black!55,thick] plot[domain=0.4:6.4,samples=70] (\x,{1.15+0.16*(\x-3.2)*(\x-3.2)});
  % 训练误差
  \draw[blue!60!black,very thick] plot[domain=0.4:6.4,samples=70] (\x,{0.35+2.6*exp(-0.45*(\x-0.4))});
  \node[font=\footnotesize,blue!60!black,anchor=west] at (6.5,0.5) {训练误差};
  % 验证误差（带噪声）
  \draw[orange!80!black,very thick] plot[smooth] coordinates {(0.4,2.48) (0.9,1.92)
    (1.4,1.74) (1.9,1.36) (2.4,1.31) (2.9,1.10) (3.4,1.20) (3.9,0.88)
    (4.4,1.44) (4.9,1.55) (5.4,2.00) (5.9,2.25) (6.4,2.85)};
  \node[font=\footnotesize,orange!80!black,anchor=west] at (6.5,2.85) {验证误差};
  % 真实风险标注
  \node[font=\footnotesize,black!70,anchor=west] at (0.35,0.42) {真实风险（看不见）};
  \draw[black!45,thin] (2.45,0.55) -- (2.88,1.08);
  % 选中的点
  \fill[red!70!black] (3.9,0.88) circle (2.6pt);
  \node[font=\footnotesize,red!70!black,anchor=north] at (3.9,0.72) {选中的点};
  \draw[<->,red!70!black] (4.05,0.90) -- (4.05,1.22);
  \node[font=\footnotesize,red!70!black,anchor=west] at (4.2,1.06) {这一截是运气};
  % 真实最优
  \fill[black!55] (3.2,1.15) circle (2.2pt);
  \draw[black!45,thin] (3.2,1.20) -- (3.2,1.36);
  \node[font=\footnotesize,black!60,anchor=south] at (3.2,1.36) {真实最优};
\end{tikzpicture}

\emph{验证曲线是真实风险加上一层噪声，取它的最小值，等于同时挑走了位置和噪声。}
\end{center}

图里三条线各说一件事。训练误差单调下降，因为模型越灵活越能贴住手上这批点，它对复杂度的选择毫无参考价值。虚线是真实风险，先降后升，最低点在偏差与方差刚好抵消的地方——这条线看不见。橙色的验证曲线是真实风险叠上一层估计噪声，你能看见的只有它。

两个位置需要分开。灰点是真实最优，红点是验证曲线的最小值。它们的横坐标常常不同，纵坐标的差距更值得留意：红点比同一横坐标上的虚线还低了一截，那一截来自``在若干个候选里取最小值''这个动作。候选越多，取最小值时捡到的负向噪声越大，红点就越往下沉。所以红点的高度不能当作这一版模型的泛化误差用——它已经被你的选择过程改造过了。

\subsection{正则化是在这条横轴上左右挪}

横轴写的是``复杂度''，但真正在动的是模型对数据波动的响应幅度，参数个数只是它的一种粗糙代理。Ridge 把系数往零收缩，Lasso 把一部分系数压成零，早停在参数还没跑远时喊停，数据增强让同一条样本以多种面貌出现，Bayesian 先验直接给参数加了一层约束，树的深度与叶子最小样本数限制了它能切多细。这些手段都不改变模型族的名义大小，却都在减小方差、增大偏差。

因此度量复杂度要用有效自由度，而不是参数个数。ridge 回归的有效自由度是

\[
\mathrm{df}(\lambda)=\trace\bigl(X(X^{\trans}X+\lambda I)^{-1}X^{\trans}\bigr)=\sum_{j}\frac{d_j^2}{d_j^2+\lambda},
\]

其中 \(d_j\) 是设计矩阵的奇异值。\(\lambda=0\) 时它等于列数，\(\lambda\to\infty\) 时趋于零，中间连续变化——同一个模型族，参数个数一个没少，有效自由度已经从 \(60\) 掉到了 \(8\)。Lasso 里数非零系数是一个方便的近似，但那些系数是看着数据选出来的，直接当自由度用会低估。

\subsection{最小验证误差不是唯一合理的落点}

交叉验证选 \(\lambda\)，做的事情就是在图里那条橙线上找最低点。做法没错，读数要打折。

先看差异是否大于噪声。开头那 \(0.802\) 与 \(0.799\)，如果各折分数的散布本身就有 \(\pm0.015\)，这个排序几乎没有信息。一个常用的做法是 one-standard-error 规则：在所有候选里，凡是误差落在最优值一个标准误以内的都算并列，从并列者中挑最简单（正则最强）的那个。它主动用一点点预测精度换取稳定性和可解释性，在特征多、样本少的场合尤其值得。前提是那个标准误得算得像样——上一节说过各折高度重叠，把 \(K\) 个折分数当独立观测算出来的标准误偏窄，用重复交叉验证的散布或逐样本 out-of-fold 损失来估更可靠。

其次要看整条曲线的形状，不只是最低点。曲线在最优点附近如果很平，说明这个旋钮在很宽的范围内都工作，选哪个都行，选宽的那一侧更安全；如果又窄又陡，说明模型对这个参数极敏感，数据稍变结论就变，这本身就是一个需要报告的事实。换个随机种子重跑一遍，看最低点跳到哪里，是最直接的稳定性检查。

\subsection{选择过程本身会过拟合验证数据}

图里红点低于虚线的那一截，是这一节最需要带走的东西。它的机制与训练误差的乐观完全一样，只是被拟合的对象从训练集换成了验证集：你在验证数据上试了 \(20\) 个 \(\lambda\)，就相当于用它训练了一个只有 \(20\) 个候选的模型；试了 \(500\) 组超参数、又比了 \(6\) 种特征方案、还调了三次决策阈值，验证集就被拟合得相当彻底了。同一份数据被查看的次数越多，它给出的分数就越不像未见数据上的表现。

后果不止在于那个分数偏高。选出来的模型的系数、普通标准误和 \(p\) 值同样过于乐观——它们是按照``这个模型是事先指定的''算出来的，而它并不是。用最终选中的模型重新拟合一遍再当作预先指定的模型来做推断，是这类错误里最常见的一种。

补救办法只有一类：让最后的评估落在没参与过选择的数据上。手段有三种。独立测试集最简单，代价是它只能用一次，一旦拿它改过任何东西就退化成了验证集。嵌套交叉验证在数据紧张时更划算，外层的每一折都对``包含调参在内的整套流程''打分。第三种是把选择步骤整个塞进重抽样循环——像本单元第一节说的那样对整条流水线做 bootstrap，或用 sample splitting、选择性推断这类专门处理``选后推断''的方法，让区间把选择的不确定性也算进去。

报告的时候，把这几件事分开写清楚：候选是怎么产生的，超参数按什么规则选的，性能在哪份数据上量的，最终模型是否在全量上重训过，测试集是不是真的只用了一次。这几行字比多报两位小数有用得多。

\subsection{信息准则走的是另一条路}

如果模型有明确的似然，还有一条不需要重抽样的近路。AIC 用

\[
\mathrm{AIC}=-2\ell(\hat\theta)+2k
\]

在拟合优度上扣掉一个复杂度惩罚，其中的 \(2k\) 正是``用同一份数据拟合又评价''所产生的乐观的一阶修正，因而它瞄准的是期望的样本外 KL 风险。BIC 把惩罚换成 \(k\log n\)，随样本量增长而变重，在候选集合包含真模型等条件下具有选择一致性。两者目标不同：前者偏向预测效率，后者偏向从有限候选里认出一个真模型；惩罚更重不等于更严谨。

用它们要守住几条边界。比较只能在同一份数据、同一种似然常数约定下进行；正则化或平滑模型里的 \(k\) 要换成有效自由度；小样本要用 AICc 修正；候选模型全都错、或维数与样本量同阶时，两者的经典解释都会失效。交叉验证不需要似然、可以直接对准你真正关心的损失，代价是计算量和估计方差，并且对切分方式敏感。在一些正则模型与损失下留一交叉验证与 AIC 渐近相通，但这不构成普遍等价。

\subsection{U 形不是普遍定律}

最后提一句边界。上面那条先降后升的曲线是经典统计的典型图景，在参数量远超样本量的现代模型上会变形：误差在插值点附近冲高之后可能再次下降，出现所谓的双下降。此时``参数多就一定方差大''不再成立，起作用的是优化算法与参数化方式共同带来的隐式正则。图里的定性结论——训练误差不能用于选择、验证读数会被选择过程污染——照旧适用；而``把复杂度调到中间''这条操作建议，需要换成对实际曲线的测量。

这一单元到此走完一条线：bootstrap 把抽样分布重放出来，置换检验把随机化机制重放出来，交叉验证把训练流程重放出来，而模型选择是在这些重放的读数上做决定，因此必须给决定本身再留一份未见数据。下一单元换一个位置提问——模型已经给出了概率，接下来该采取什么行动，错误的代价怎样进入这个决定，见\hyperref[ch:092]{``统计决策与可信预测：模型给出概率之后怎样行动''}。
